\documentclass[11pt,a4paper]{article}
\usepackage[utf8]{inputenc}
\usepackage[T1]{fontenc}
\usepackage[polish]{babel}
\usepackage{lmodern}
\usepackage{geometry}
\geometry{margin=27mm}
\usepackage[hidelinks]{hyperref}

\title{Synergia EEG, mimiki i zachowania w pomiarze powrotu do regulacji}
\author{}
\date{Koncepcja badawcza, IDEA / DEFER, wersja 0.1.2}

\begin{document}
\maketitle

\begin{abstract}
Praca przedstawia odroczoną koncepcję oceny, czy łączne wykorzystanie EEG, mimiki i telemetrii zachowania poprawia pomiar powrotu do regulacji (ang.~\emph{recovery}) w porównaniu z poszczególnymi modalnościami. Nie ustalono jeszcze definicji zdarzenia, wyniku pierwszorzędowego ani estymandu. Tekst nie opisuje przeprowadzonego badania i nie zawiera wyników empirycznych.
\end{abstract}

\section{Wprowadzenie}
Powrót do regulacji może obejmować odmienne zjawiska: ponowne podjęcie zadania, zmianę zachowania obserwowalnego i zmiany sygnałów fizjologicznych. Zgodność tych wskaźników nie może być zakładana z góry. Przed rozpoczęciem badania należy ustalić, który konstrukt jest przedmiotem pomiaru, a które sygnały pełnią jedynie funkcję predyktorów.

\section{Cel i pytania badawcze}
Planowane pytanie dotyczy tego, czy fuzja EEG, mimiki i telemetrii zadania wnosi wartość predykcyjną ponad model oparty na samym zachowaniu. Rozstrzygnięcia wymaga operacyjna definicja początku i końca epizodu powrotu do regulacji, okno obserwacji, jednostka analizy oraz miara jakości predykcji. Nie sformułowano jeszcze potwierdzonych hipotez kierunkowych.

\section{Materiał i metody}
Rozważany projekt zakłada rejestrację modalności w tych samych epizodach oraz synchronizację ich znaczników czasu. Ocena jakości sygnału i dopuszczalnego opóźnienia między kanałami powinna poprzedzać analizę. Wariant referencyjny oparty wyłącznie na telemetrii zadania należy porównać z modelami rozszerzanymi o mimikę, EEG i obie modalności. Należy oddzielić cechy dostępne w chwili predykcji od informacji o późniejszym wyniku, aby uniknąć przecieku czasowego.

Badanie pozostaje odroczone (\texttt{DEFER}). Nie potwierdzono kompletnego protokołu, gotowości aparatury, uprawnień do danych, procedur etycznych ani pozyskania próby. Propozycje stanowiska i oprogramowania opisane w \texttt{study.md} nie stanowią dowodu ich wdrożenia.

\section{Plan analizy}
Po zdefiniowaniu wyniku i estymandu należy z góry wskazać metryki jakości i kalibracji, scenariusz brakujących modalności oraz analizę wrażliwości na synchronizację. Wstępnym wariantem walidacji jest podział z wyłączeniem całego uczestnika (\emph{leave-one-participant-out}); wszystkie epizody jednego uczestnika muszą należeć wyłącznie do jednej części podziału. Dobór cech, strojenie hiperparametrów i kalibracja powinny odbywać się wewnątrz danych treningowych. Metryka i sposób agregacji na poziomie uczestnika wymagają uprzedniego ustalenia. Analiza nie została wykonana.

\section{Wyniki}
Brak danych empirycznych i wyników analiz. Nie należy interpretować projektu porównań modeli jako dowodu synergii.

\section{Dyskusja i ograniczenia}
Potencjalne ograniczenia obejmują artefakty EEG, niejednoznaczność mimiki jako wskaźnika stanu afektywnego, nierówną jakość kanałów, brakujące obserwacje i różnice między uczestnikami. Dodatkowa trafność predykcyjna nie dowodzi przyczynowej roli modalności ani użyteczności klinicznej. Odrębność planowanej pracy wobec MSC-ADHD-P02 oraz MSC-ADHD-P05 wymaga oceny na podstawie rzeczywistych protokołów.

\section{Wnioski}
Na etapie \texttt{IDEA / DEFER} możliwe jest jedynie określenie warunków sprawdzalności koncepcji. Rozpoczęcie badań wymaga potwierdzenia kanonicznego projektu, protokołu, praw do danych i niezbędnych zgód.

\section*{Bibliografia}
Pozycje bibliograficzne zostaną dodane dopiero po weryfikacji źródeł i ich związku z konkretnymi twierdzeniami.

\end{document}
